\section{Bonding, Lewis Structures, Formal Charge, and VSEPR}

\subsection{Ctrl+F keywords: Lewis, formal charge, resonance, VSEPR, bond angle, planar, polarity, dipole, hybridization}
Use for drawing or checking Lewis structures, choosing resonance forms, molecular geometry, bond angles, polarity, and planar/non-planar questions.

\subsection{Symbols and notation}
\referencetable{tab:bonding:symbols}{Bonding, Lewis, and VSEPR symbols.}
\begin{tabular}{@{}>{$}l<{$} p{10.8cm}@{}}
\toprule
\text{Symbol} & \text{Meaning and usage}\\
\midrule
N_{\mathrm{val}},\ q & Total number of valence electrons and overall ionic charge.\\
\mathrm{FC},\ V & Formal charge and valence-electron count of the selected atom.\\
N_{\mathrm{nonbonding}},\ N_{\mathrm{bonding}} & Nonbonding and bonding electrons assigned in a Lewis structure.\\
\mathrm{SN} & Steric number: bonded atoms plus lone pairs on the central atom.\\
\vec{\mu}_i,\ \vec{\mu}_{\mathrm{mol}},\ i & Individual bond dipole, vector sum for molecular dipole, and bond index in the sum.\\
\chi,\ \Delta\chi & Electronegativity and electronegativity difference.\\
N_{\mathrm{bonding}},\ N_{\mathrm{antibonding}} & Electrons in bonding and antibonding molecular orbitals in the bond-order formula.\\
\bottomrule
\end{tabular}

\subsection{Lewis structures}

\subsubsection{Valence electron count}
\begin{equation}
\label{eq:bonding:valence-electrons}
N_{\mathrm{val}}=\sum_i N_{\mathrm{val},i}-q
\end{equation}
For an ion, subtract positive charge and add negative charge. Total Lewis electrons must equal \(N_{\mathrm{val}}\).

\subsubsection{Formal charge}
\begin{equation}
\label{eq:bonding:formal-charge}
\mathrm{FC}=V-\left(N_{\mathrm{nonbonding}}+\frac12N_{\mathrm{bonding}}\right)
\end{equation}
Preferred structures usually minimize formal charges and place negative formal charge on more electronegative atoms.

\subsubsection{Resonance structures}
\begin{equation}
\label{eq:bonding:resonance-rule}
\text{resonance forms differ in electron placement, not atom placement}
\end{equation}
Equivalent resonance forms indicate delocalization and equalized bond lengths. Do not move nuclei when drawing resonance.

\subsubsection{Octet, expanded octet, electron-deficient exceptions}
\begin{equation}
\label{eq:bonding:octet-rule}
\text{period 2 atoms normally obey octet}, \qquad
\text{period }3+\text{ may expand octet}
\end{equation}
Hydrogen has a duet. Boron can be electron deficient; sulfur, phosphorus, iodine, and xenon can appear hypervalent in course examples.

\subsection{Exam recipe: Lewis structure and formal charge}
\textbf{Use when:} the problem asks for a valid Lewis structure, resonance contributors, formal charges, nitrate, or oxalate.

\textbf{Given / Find:} atom identities and charge \(q\); find electron placement, resonance forms, and formal charges.

\textbf{Required references:} Equation~\ref{eq:bonding:valence-electrons}, Equation~\ref{eq:bonding:formal-charge}, Equation~\ref{eq:bonding:resonance-rule}, and Equation~\ref{eq:bonding:octet-rule}.

\textbf{Steps:}
\begin{enumerate}
    \item Compute \(N_{\mathrm{val}}\) from Equation~\ref{eq:bonding:valence-electrons}; reserve exactly that many electrons.
    \item Choose a skeleton: hydrogen is terminal; otherwise put the least electronegative suitable atom centrally.
    \item Draw single bonds, complete terminal-atom octets or hydrogen duets, then place remaining electrons centrally.
    \item If a second-period central atom is electron deficient, convert adjacent lone pairs to multiple bonds without exceeding its octet.
    \item Calculate \(\mathrm{FC}\) on every atom using Equation~\ref{eq:bonding:formal-charge}; verify that the sum is \(q\).
    \item Draw equivalent electron-placement alternatives as resonance forms and select contributors with minimized, chemically sensible formal charges.
\end{enumerate}

\textbf{Checks:} electrons drawn equal \(N_{\mathrm{val}}\); formal charges sum to \(q\); no second-period atom exceeds an octet.

\textbf{Common traps:} moving atoms between resonance forms; ignoring ionic charge in the count; expanding a second-period octet.

\subsection{VSEPR geometry}

\subsubsection{Steric number and electron domains}
\begin{equation}
\label{eq:bonding:steric-number}
\mathrm{SN}=\#\text{bonded atoms}+\#\text{lone pairs on central atom}
\end{equation}
Multiple bonds count as one electron domain in VSEPR. Lone pairs compress bond angles.

\subsubsection{Common VSEPR geometries and bond angles}
\referencetable{tab:bonding:vsepr-geometries}{Common central-atom VSEPR geometries and ideal or approximate bond angles.}
\begin{tabular}{@{}cccc@{}}
\toprule
\(\mathrm{SN}\) & Lone pairs & Molecular geometry & Bond angle\\
\midrule
\(2\) & \(0\) & linear & \(180^\circ\)\\
\(3\) & \(0\) & trigonal planar & \(120^\circ\)\\
\(4\) & \(0\) & tetrahedral & \(109.5^\circ\)\\
\(4\) & \(1\) & trigonal pyramidal & \(\approx107^\circ\)\\
\(4\) & \(2\) & bent & \(\approx105^\circ\)\\
\(5\) & \(0\) & trigonal bipyramidal & \(90^\circ,120^\circ\)\\
\(6\) & \(0\) & octahedral & \(90^\circ\)\\
\(6\) & \(1\) & square pyramidal & \(\approx90^\circ\)\\
\(6\) & \(2\) & square planar & \(90^\circ\)\\
\bottomrule
\end{tabular}
Exam wording often asks ``planar'', ``non-planar'', ``bond angle'', ``pyramidal'', or ``bent''.

\subsubsection{Planar molecule recognition}
\begin{equation}
\label{eq:bonding:planarity-rule}
\text{linear, trigonal planar, square planar, and many aromatic rings are planar}
\end{equation}
Tetrahedral and trigonal pyramidal centers are non-planar around that atom. A molecule is planar only when all atoms can lie in one plane.

\subsubsection{Molecular polarity}
\begin{equation}
\label{eq:bonding:molecular-dipole}
\vec{\mu}_{\mathrm{mol}}=\sum_i\vec{\mu}_i
\end{equation}
Polar bonds can cancel in symmetric geometries. Lone pairs often prevent cancellation and make the molecule polar.

\subsection{Exam recipe: molecular geometry, bond angle, and planarity}
\textbf{Use when:} the problem asks for VSEPR shape, approximate bond angle, planar/non-planar, or central-atom geometry.

\textbf{Given / Find:} a molecular formula or Lewis structure; find geometry, requested angle, and planarity.

\textbf{Required references:} Equation~\ref{eq:bonding:steric-number}, Table~\ref{tab:bonding:vsepr-geometries}, and Equation~\ref{eq:bonding:planarity-rule}.

\textbf{Steps:}
\begin{enumerate}
    \item Draw or verify the Lewis structure, including central-atom lone pairs.
    \item Count bonded atoms and lone pairs around each relevant center; compute \(\mathrm{SN}\) by Equation~\ref{eq:bonding:steric-number}.
    \item Match \(\mathrm{SN}\) and lone-pair count to Table~\ref{tab:bonding:vsepr-geometries}.
    \item Report the listed angle; state that lone pairs compress an ideal angle where applicable.
    \item Test planarity from the local geometries: one tetrahedral or trigonal-pyramidal arrangement containing required atoms makes the molecule non-planar.
\end{enumerate}

\textbf{Checks:} multiple bonds were counted as one domain; the stated geometry matches the lone-pair count; the planarity claim concerns all atoms requested.

\textbf{Common traps:} counting bond order as domains; reporting electron-domain geometry instead of molecular geometry; calling a molecule planar from one planar fragment.

\subsection{Exam recipe: molecular polarity}
\textbf{Use when:} the problem asks for polar/nonpolar, dipole moment, bond-dipole cancellation, or polarity ranking.

\textbf{Given / Find:} bond identities and molecular structure; find whether \(\vec{\mu}_{\mathrm{mol}}\) is zero or nonzero.

\textbf{Required references:} Equation~\ref{eq:bonding:molecular-dipole}, Equation~\ref{eq:bonding:bond-polarity}, and Table~\ref{tab:bonding:vsepr-geometries}.

\textbf{Steps:}
\begin{enumerate}
    \item Identify polar bonds by electronegativity difference using Equation~\ref{eq:bonding:bond-polarity}.
    \item Obtain the three-dimensional molecular geometry from the Lewis structure and VSEPR table.
    \item Draw bond-dipole directions toward more electronegative atoms; include asymmetry caused by substituent changes or lone pairs.
    \item Add vectors by symmetry using Equation~\ref{eq:bonding:molecular-dipole}: complete cancellation means nonpolar; a residual vector means polar.
\end{enumerate}

\textbf{Checks:} a molecule with no polar bonds is nonpolar; equal symmetrically arranged dipoles cancel; an unsymmetrical arrangement usually does not.

\textbf{Common traps:} deciding molecular polarity from a single bond; ignoring shape or lone pairs; assuming every symmetric-looking drawing is three-dimensionally symmetric.

\subsection{Bonding formulas}

\subsubsection{Bond polarity from electronegativity}
\begin{equation}
\label{eq:bonding:bond-polarity}
\Delta\chi=|\chi_A-\chi_B|
\end{equation}
Larger electronegativity difference gives a more polar bond. Geometry decides whether the whole molecule is polar.

\subsection{Exam recipe: ionic versus covalent bond character}
\textbf{Use when:} the task asks which bond has greater ionic character, greater covalent character, or the smallest/largest polarity.

\textbf{Given / Find:} Given several bonded atom pairs and electronegativity values or periodic positions; rank bond character.

\textbf{Required references:} Equation~\eqref{eq:bonding:bond-polarity}; use supplied electronegativity values when provided.

\textbf{Steps:}
\begin{enumerate}
    \item For each bonded pair, obtain \(\chi_A\) and \(\chi_B\) from the provided table or accepted periodic trend data.
    \item Calculate \(\Delta\chi=|\chi_A-\chi_B|\) for every pair using Equation~\eqref{eq:bonding:bond-polarity}.
    \item Rank increasing ionic character by increasing \(\Delta\chi\); reverse this ranking for increasing covalent character.
    \item For a bond between identical atoms, set \(\Delta\chi=0\) and classify it as the most purely covalent option in that comparison.
\end{enumerate}

\textbf{Checks:} The ranking is based on a bond, not whole-molecule polarity; larger \(\Delta\chi\) cannot be labelled more covalent in the same set.

\textbf{Common traps:} using molecular shape to rank bond character; reversing the covalent ranking; overlooking that a homonuclear bond has zero electronegativity difference.

\subsubsection{Molecular orbital bond order}
\begin{equation}
\label{eq:bonding:mo-bond-order}
\text{bond order}=\frac{N_{\mathrm{bonding}}-N_{\mathrm{antibonding}}}{2}
\end{equation}
Higher bond order usually means shorter, stronger bonds. Zero bond order predicts no stable bond in simple MO treatment.
